\subsection{KG Query Engine}

The query engine supports a SPARQL-inspired syntax with four execution strategies.

\subsubsection{SPARQL Syntax}

Supported query structure:

\begin{verbatim}
SELECT ?var1 ?var2 ...
WHERE {
    ?s ?p ?o .
    FILTER(condition)
    OPTIONAL { ... }
}
ORDER BY ?var ASC|DESC
LIMIT n
OFFSET m
\end{verbatim}

\subsubsection{Query Planning}

Query execution follows a two-phase planning approach:

\textbf{Logical Phase}: The query is translated into a tree of logical operators:
\begin{itemize}
    \item \texttt{LogicalScan}: Full sequential scan of the triple table
    \item \texttt{LogicalIndexSeek}: Index-based access by subject, predicate, or object
    \item \texttt{LogicalFilter}: Predicate evaluation on rows
    \item \texttt{LogicalJoin}: Combination of two sub-plans (nested-loop or hash join)
    \item \texttt{LogicalProject}: Column selection
    \item \texttt{LogicalSort}: Row ordering
    \item \texttt{LogicalLimit}: Row count restriction
    \item \texttt{LogicalAggregate}: Grouping and aggregation
\end{itemize}

\textbf{Physical Phase}: Each logical operator is converted to a physical execution strategy.

\subsubsection{Optimization}

The cost-based optimizer applies three transformations:

\begin{enumerate}
    \item \textbf{Filter pushdown}: Moves filter predicates closer to data sources to reduce intermediate row counts early.
    \item \textbf{Join reordering}: Reorders joins so that smaller intermediate results are processed first, reducing overall cost.
    \item \textbf{Index selection}: Chooses between SPO, POS, and OPS indexes based on which columns are bound in the query pattern.
\end{enumerate}

Cost estimation uses a simple model:
\[
\text{cost} = N_{\text{pages}} \times C_{\text{seq}} + N_{\text{rows}} \times C_{\text{cpu}}
\]

\subsubsection{Execution}

The physical plan is executed by a tree-walking interpreter that materializes intermediate results. For large result sets, streaming execution can be used to bound memory usage.
